\chapter{Literature Survey}
\label{ch:literature-survey}

\section{Overview}
Zero-day detection requires identifying previously unseen attacks without relying
on signature databases. Recent research emphasizes learning normal behavior from
structure-rich data and detecting deviations using self-supervised objectives. The
papers chosen for this project align with three pillars: (i) self-supervised graph
representation learning for intrusion detection, (ii) zero-day detection using
machine learning without prior attack labels, and (iii) continual learning and
alert correlation to make the system practical over time. This survey correlates
each paper directly with the modules implemented in the project.

\section{Self-Supervised Graph Learning for Intrusion Detection}
Guerra et al. present a self-supervised method for learning graph representations
for intrusion detection, demonstrating that graph structure alone can provide
strong anomaly signals without labeled attacks \cite{guerra2025graphids}. This
directly motivates the use of graph autoencoders in our system, where behavioral
graphs are reconstructed and high reconstruction error indicates an anomaly.

Hayat et al. study self-supervised time-aware heterogeneous hypergraph learning
for dynamic graph-level classification \cite{hayat2025self}. Their work supports
two key design choices in our project: (i) representing heterogeneous entities
(process, file, socket) and (ii) capturing time-evolving behavior via windows of
events. Our graph construction module follows this insight by building time-window
behavior graphs and preserving type-specific interactions.

\section{Zero-Day Detection Using Machine Learning}
Dai et al. propose a zero-day intrusion detection model for unseen data using
machine learning \cite{dai2025zero}. This work reinforces the need for detection
without prior attack labels, aligning with our self-supervised training strategy.
In our pipeline, the model is trained only on benign behavior, and anomalies are
detected as deviations in reconstruction error rather than signature matches.

\section{Alert Correlation and Practical Detection Pipelines}
Lanigan et al. focus on alert correlation for intelligent threat detection and
response \cite{lanigan2025alert}. Their study highlights that detection systems
must produce interpretable, correlatable alerts rather than isolated anomaly
scores. This informs our alerting and logging module, where each anomaly score is
converted into structured alerts with severity for downstream correlation.

\section{Continual Learning and Model Adaptation}
Continual learning is critical because normal behavior shifts over time. Guo et
al. present continual learning strategies for intrusion detection under evolving
threats \cite{guo2025continual}, while Prasath et al. analyze multiple continual
learning models in IDS settings \cite{prasath2025analysis}. Lin et al. further
emphasize dynamic model updates driven by cyber threat intelligence \cite{lin2025evolving}.
Together, these works justify the project's modular architecture and the planned
future extension where the model can be updated with new normal behavior without
forgetting prior patterns.

\section{Module-to-Literature Mapping}
Based on the above studies, each module in the project is grounded in recent
research:
\begin{itemize}
    \item Graph autoencoder learning and reconstruction-based detection are
    supported by self-supervised graph IDS research \cite{guerra2025graphids}.
    \item Heterogeneous, time-windowed graph construction aligns with time-aware
    dynamic graph learning \cite{hayat2025self}.
    \item Zero-day, label-free detection motivation is reinforced by recent
    machine-learning-based IDS work \cite{dai2025zero}.
    \item Alert structuring and interpretability are informed by alert
    correlation research \cite{lanigan2025alert}.
    \item Long-term adaptability and model updates are grounded in continual
    learning IDS studies \cite{guo2025continual,prasath2025analysis,lin2025evolving}.
\end{itemize}

\section{Summary}
The selected literature provides a coherent foundation for the proposed system.
Self-supervised graph learning enables anomaly detection without labels, dynamic
graph modeling supports heterogeneous behavioral data, and continual learning
ensures long-term adaptability. By aligning each system module with the cited
studies, the project follows current research directions while addressing the
practical requirements of zero-day intrusion detection.
